\section*{Цель}
Целью выпускной квалификационной работы является проектирование и разработка платформы TechEventsHub для автоматизированного сбора, интеллектуального анализа и персонализированной доставки анонсов профильных IT-мероприятий из внешних неструктурированных источников (Telegram-каналов и Web-страниц).

Данная платформа станет эффективным инструментом, ускоряющим и упрощающим процесс поиска релевантных событий, позволит снизить информационную перегрузку IT-специалистов и будет способствовать повышению посещаемости IT-мероприятий.

\section*{Задачи}
Задачи проекта разделяются на две части: исследовательскую и программную, а именно:

\begin{itemize}
  \item{Исследовательская часть}
    \begin{itemize}
        
    \item Провести опрос целевой аудитории IT-специалистов для изучения текущего пользовательского опыта, выявления основных барьеров при поиске анонсов профильных мероприятий и определения актуальных потребностей пользователей.

    \item Провести сравнительный анализ существующих на рынке решений (платформ-агрегаторов IT-мероприятий), выявленных в ходе исследования, для определения их функциональных ограничений и недостатков.

    \item На основе полученных аналитических данных сформировать комплекс функциональных и нефункциональных требований к разрабатываемой платформе, а также обосновать выбор программного и технологического инструментария.

    \end{itemize}
  
  \item{Программная часть}
    \begin{itemize}
        
    \item Спроектировать отказоустойчивую микросервисную архитектуру системы и реляционную схему базы данных для хранения пользовательских настроек и метаданных.
    
    \item Разработать подсистему фонового сбора данных (скрейпинга) из открытых Telegram-каналов и Web-страниц с механизмами валидации источников.
    
    \item Обеспечить бесперебойное межсервисное взаимодействие с использованием брокера сообщений (Apache Kafka) и кэширования (Redis).
    
    \item Интегрировать методы обработки естественного языка на базе LLM YandexGPT для извлечения структурированных данных из сырых текстов анонсов.
    
    \item Разработать механизм строгой типизированной фильтрации и подбора событий на основе индивидуальных предпочтений пользователей.
    
    \item Спроектировать подсистему единой идентификации пользователей с поддержкой Yandex OAuth и механизмом глубокого связывания Web-аккаунтов с Telegram-профилями.
    
    \item Разработать маршрутизатор омниканальных уведомлений для доставки сформированных дайджестов пользователям в Telegram и на E-mail.
    
    \item Реализовать клиентское Web-приложение (на базе Next.js) для предоставления интерактивного графического интерфейса управления подписками и дашборда анонсов событий.
    
    \item Настроить инфраструктуру наблюдаемости (Prometheus, Grafana) и сбор RED-метрик сервисов. 
    
\end{itemize}

\end{itemize}


\section*{Актуальность работы}

Поскольку в наши дни информационные технологии стремительно развиваются, для продуктивной работы IT-специалистам необходимо непрерывное обучение, обмен опытом и расширение профессиональных связей. 
Согласно современным исследованиям и производственной практике, посещение профильных мероприятий (конференций, митапов, воркшопов) является неотъемлемой частью профессионального развития и карьерного роста.

В аналитических материалах компании Momentive Software~\cite{attending_conferences_benefits} отмечается, что участие в профильных конференциях решает сразу несколько важных задач для специалиста. 
Во-первых, оно предоставляет возможность очного общения с другими специалистами индустрии, которое создает более прочный фундамент для дальнейшего сотрудничества, чем цифровые коммуникации. 
Во-вторых, поскольку тематические профессиональные мероприятия посещают специалисты из различных смежных областей и находящиеся на разных этапах карьеры, общение с ними позволяет получить актуальные пприкладные знания и понимание технологических трендов. 
Наконец, активное участие в мероприятиях повышает профессиональную видимость специалиста и способствует формированию его статуса как эксперта в отрасли~\cite{attending_conferences_benefits}.

Высокая значимость участия в профессиональных мероприятиях также подтверждается академическими исследованиями. 
В частности, в исследовании К. Хаусса~\cite{hauss2020} доказано, что очные конференции представляют собой значимые социальные события, оказывающие влияние как на развитие отдельных специалистов, так и профессиональных сообществ.
Исследование показывает, что для более 80\% участников очное присутствие на профильных событиях приводит к прямому обмену методами и результатами работы, а для значительной части (до 38\%) - становится отправной точкой для новых совместных проектов и исследований. 
Кроме того, автор подчеркивает, что именно на таких мероприятиях происходит обмен передовым, <<еще не опубликованными>> знаниями и неформальной карьерной информацией, доступ к которой невозможен через онлайн каналы коммуникации~\cite{hauss2020}.

Тем не менее, несмотря на продемонстрированную высокую ценность посещения профессиональных мероприятий, процесс их поиска сопряжен с рядом проблем, что подтверждает актуальность разработки специализированной платформы для агрегации и систематизации анонсов событий.

\subsection*{Особенности отрасли}
Сегодня подход к поиску IT-мероприятий изменяется: аудитория уходит от массовых агрегаторов к разрозненным, более узкоспециализированным источникам. 
Большое количество профильных и локальных событий сегодня анонсируется только в тематических Telegram-каналах и на независимых сайтах организаций. 
Подобная разрозненность источников приводит к следующим отраслевым проблемам:
\begin{itemize}
\item \textbf{Информационный шум.} Пользователи вынуждены подписываться на десятки каналов, где анонсы мероприятий могут смешиваться с нерелевантными новостями и рекламой.
\item \textbf{Высокая когнитивная нагрузка.} Отсутствие единого формата постов заставляет пользователей тратить время на вычитывание базовой информации (дата, место, стоимость события).
\item \textbf{Боязнь пропустить интересное (Fear of missing out)}~\cite{fomo_def}. Из-за лент, основанных на рекомендательных алгоритмах, и большого потока сообщений IT-специалисты могут пропускать регистрацию на бесплатные или популярные мероприятия, где билеты заканчиваются за несколько часов.
\end{itemize}
Таким образом, ручной поиск актуальных мероприятий для посещения становится неэффективным, что снижает конверсию из просмотра анонса в посещение мероприятия.

\subsection*{Ценностные предложения}
Платформа TechEventsHub предлагает пользователям подход осознанного управления своим информационным полем. 
Главная ценность продукта заключается в следующем:
\begin{itemize}
\item \textbf{Экономия времени.} Сокращение времени с более 30 минут <<ручного>> мониторинга источников в неделю до 5 минут просмотра структурированного дайджеста.
\item \textbf{Снижение когнитивной нагрузки.} Использование LLM YandexGPT для извлечения сути и основных параметров события из неструктурированных текстов анонсов и формирование единообразных карточек анонсов (Повестка мероприятия / Место / Время / Цена / Ссылка на источник).
\item \textbf{Персонализация контента.} Типизированная фильтрация (по датам, технологическому стеку, городу и стоимости) позволяет предоставить пользователю релевантные анонсы из выбранных им источников в соотвествии с настроенными им параметрми поиска.
\item \textbf{Улучшение пользовательского опыта.} Наличие Web-дашборда для удобной настройки фильтров в связке с омниканальной доставкой дайджестов анонсов (в Telegram-боте и по E-mail) и генерацией \texttt{.ics}-файлов для быстрого добавления понравившихся событий в электронный календарь.
\end{itemize}

\subsection*{Рисковые гипотезы}
В рамках проектирования и реализации платформы выдвигается ряд рисковых гипотез:
\begin{enumerate}
\item \textbf{Качество извлечения и суммаризации данных ИИ}
  \begin{itemize}
    \item Предполагается, что современные LLM способны с точностью более 95\% извлекать сложные метаданные (дату, цену, формат) из текстов анонсов. 
    Риск <<галлюцинаций>> модели планирается митигировать промпт-инжинирингом и программной валидацией результатов суммаризации.
  \end{itemize}

\item \textbf{Поведенческая гипотеза} 
  \begin{itemize}
    \item Предполагается, что представители целевой аудитории\\(IT-специалисты) предпочтут самостоятельно добавлять ссылки на источники анонсов вместо использования классических рекомендательных лент, что обеспечит высокий показатель возврата пользователей платформы.
  \end{itemize}

\item \textbf{Технические ограничения}
    \begin{itemize}
        \item Существует риск блокировок со стороны Telegram MTProto API из-за автоматизированного сбора данных. 
        Данный риск планируется митигировать использованием пагинации, эмуляции задержек и реализации паттерна Circuit Breaker для обеспечения стабильной работы сервиса скрейпинга.
        \item Возможны сложности с парсингом Web-сайтов из-за ограничений их собственников.
    \end{itemize}
\end{enumerate}

\vspace{1em}

Таким образом, можно считать, что актуальность работы продемонстрирована.
